\documentclass[12pt]{article}
\usepackage[utf8]{inputenc}
\usepackage[T1]{fontenc}
\usepackage{hyperref}
\title{Ancient Ligurian Language Resources for K-12 Curriculum Development}
\author{Compiled by Claude Sonnet 5.5 (claude-sonnet-5-5) for Andr\'as Kornai}
\date{}
\begin{document}
\maketitle

\section{Introduction}
Ancient Ligurian is barely attested. There are no continuous texts. The evidence consists of place, river, tribal and personal names, a few glosses in Greek and Latin authors, and a few disputed inscriptions. The most important document is the Polcevera Tablet (117 BCE), a Roman arbitration record that preserves indigenous place names. The Lunigiana statue-stelae carry inscriptions in an Etruscan-derived alphabet. Their attribution to Ligurian and their transcription and interpretation are disputed. The language is generally thought to be Indo-European, but its position is debated: some scholars treat it as Celtic, others as a separate branch. No educational materials in or about the language were found.

\section{Primary Educational Resources}
No primers, textbooks or school programmes were found.

\section{Linguistic and Reference Materials}
\begin{itemize}
\item Antigone Journal, ``What Did Palaeo-European Peoples Write? Local Languages of the Western Mediterranean'': an accessible overview that includes the western Mediterranean languages. \url{https://antigonejournal.com/2023/05/western-mediterranean-ancient-languages/}
\item Wikipedia, ``Liguri'': background on the Ligures as a people. \url{https://en.wikipedia.com/wiki/Liguri}
\item Bernard Mees, academia.edu papers page: a list of an academic's papers. I did not verify which papers concern Ligurian. \url{https://monash.academia.edu/Bernard Mees/Papers}
\end{itemize}

\section{Religious and Community Resources}
None were found for the ancient language.

\section{Digital and Audio Resources}
No audio, software or corpus resources dedicated to Ligurian were found.

\section{Curriculum Development Context}
There is nothing to teach as a language beyond onomastic material. A heritage unit could present the Ligures, the Polcevera Tablet place names and the Lunigiana stelae as local history. It would need to say plainly that the linguistic evidence is thin and contested.

\end{document}